\documentclass[12pt]{article}
\usepackage[a4paper,top=2.8cm,bottom=0.8cm,left=2.5cm,right=2.5cm]{geometry}
\usepackage{graphicx}
\usepackage{mathptmx} 
\usepackage{titlesec}
\usepackage{xcolor}
\usepackage{amsmath}
\usepackage{amsmath, amssymb}
\usepackage{enumitem}
\usepackage{setspace}
\usepackage{titling}
\usepackage{lipsum}
\usepackage{tikz}
\usepackage{changepage}
\usepackage{float}
\usepackage{caption}
\usepackage{subcaption}
\usepackage{fancyhdr}
\pagestyle{empty}
\setlength{\parindent}{0pt}
\setlength{\parskip}{0.8em}
\setstretch{1.1}
\definecolor{cyanblue}{RGB}{0,153,204}
\definecolor{darkskyblue}{rgb}{0.0, 0.5, 1.0}
\definecolor{skyblue}{RGB}{135, 206, 235}
% Section styling
\titleformat{\section}{\color{cyanblue}\large\bfseries}{\thesection}{1em}{}
\titleformat{\subsection}[runin]{\color{cyanblue}\bfseries}{\thesubsection}{1em}{}[.]
\begin{document}
\pagestyle{empty} % Start with empty page style
\thispagestyle{fancy} % Apply fancy style only to the first page
\renewcommand{\headrulewidth}{0pt} % Remove header rule
\fancyhead[L]{% Left header
        \includegraphics[width=8cm, height=1.7cm]{img3.png} 
        }
\fancyhead[R]{% Right header
    Name: A.Likhitha \\
    Batch: COMETFWC019 \\
    Date: 17 may 2025
}

\vspace*{0.4cm}
\noindent
{\color{cyanblue}\makebox[\textwidth]{\text{Quadratic Equations} \hfill \text{39}}\\[-0.6em]
\rule{\textwidth}{3pt}}
\noindent
\smallskip
Sridharacharya (C.E. 1025) derived a formula, now known as the quadratic formula, (as quoted by Bhaskara II) for solving a quadratic equation by the method of completing the square. An Arab mathematician Al-Khwarizmi (about C.E. 800) also studied quadratic equations of different types. Abraham bar Hiyya Ha-Nasi, in his book ‘Libre embadorum’ published in Europe in C.E. 1145 gave complete solutions of different quadratic equations.

\hspace{3em}In this chapter, you will study quadratic equations, and various ways of finding\\ their roots. You will also see some applications of quadratic equations in daily life\\ situations.\\\\
{\color{cyanblue}\bfseries 4.2 \quad Quadratic Equations}\\
\smallskip
A quadratic equation in the variable $x$ is an equation of the form $ax^2 + bx + c = 0$, where $a, b, c$ are real numbers, $a \ne 0$. For example, $2x^2 + x - 300 = 0$ is a quadratic equation. Similarly, $2x^2 - 3x + 1 = 0$, $4x - 3x^2 + 2 = 0$ and $1 - x^2 + 300 = 0$ are also quadratic equations.

\hspace{3em}In fact, any equation of the form $p(x) = 0$, where $p(x)$ is a polynomial of degree 2, is a quadratic equation. But when we write the terms of $p(x)$ in descending order of their degrees, then we get the standard form of the equation. That is, $ax^2 + bx + c = 0$, $a \ne 0$ is called the \textbf{standard form of a quadratic equation}.

Quadratic equations arise in several situations in the world around us and in different fields of mathematics. Let us consider a few examples.

\textbf{\color{cyanblue}Example 1:} Represent the following situations mathematically:
\vspace{-1em}
\begin{enumerate}[label=(\roman*), leftmargin=0.7cm] 
\item John and Jivanti together have 45 marbles. Both of them lost 5 marbles each, and the product of the number of marbles they now have is 124. We would like to find out how many marbles they had to start with.
 
\item A cottage industry produces a certain number of toys in a day. The cost of production of each toy (in rupees) was found to be 55 minus the number of toys produced in a day. On a particular day, the total cost of production was ₹750. We would like to find out the number of toys produced on that day.
\end{enumerate}
\vspace{-1em}
\textbf{\color{cyanblue}Solution:}
\begin{enumerate}[label=(\roman*), leftmargin=0.8cm]
\vspace{-1em}
\item Let the number of marbles John had be $x$.\\
Then the number of marbles Jivanti had = $45 - x$ (Why?).\\
The number of marbles left with John, when he lost 5 marbles = $x - 5$\\
The number of marbles left with Jivanti, when she lost 5 marbles = $45 - x - 5$\\
\hspace*{11cm}$= 40 - x$
\end{enumerate}
\vspace*{4em}
\centering
Reprint 2025-2026
\newpage
\noindent
\begin{tikzpicture}
  \draw[cyanblue, line width=2pt] (0,0) -- (\linewidth,0); % Horizontal line
  \node[anchor=west, fill=white, inner sep=1pt, text=cyanblue] at (0,0.3) {\text{40}}; % Left text
  \node[anchor=east, fill=white, inner sep=1pt, text=cyanblue] at (\linewidth,0.3) {\text{\text{Mathematics}}}; % Right text
\end{tikzpicture}
\vspace{-3em}
\begin{flushleft}
\hspace{2em}{Therefore, their product = (x – 5) (40 – x)}\\
\hspace{12em}= 40x – x^2– 200 + 5x}\\
\end{flushleft}
\hspace{-5em} 
So, \quad -x^2 + 45x - 200 = 124
\hspace{6em} {\text{\text{(Given that product = 124)}}}\\
\hspace*{-22em} i.e., \quad $-x^2 + 45x - 324 = 0$\\
\hspace*{-23em} i.e., \quad $x^2 - 45x + 324 = 0$\\
\hspace{-4em}
Therefore, the number of marbles John had, satisfies the quadratic equation
\[
{x^2 - 45x + 324 = 0}
\]
\hspace{-7em}
which is the required representation of the problem mathematically.\\
\vspace{0.4em}
\hspace{-13em}
(ii) Let the number of toys produced on that day be $x$.\\
\vspace{0.4em}
\hspace*{-4em} Therefore, the cost of production (in rupees) of each toy that day $= 55 - x$\\
\vspace{0.4em}
\hspace*{-8em} So, the total cost of production (in rupees) that day $= x(55 - x)$\\
\vspace{0.4em}
\hspace*{-16em} Therefore,\hspace{6em}\quad $x(55 - x) = 750$\\
\vspace{0.4em}
\hspace*{-16em} i.e.,\hspace{8em} \quad $55x - x^2 = 750$\\
\vspace{0.4em}
\hspace*{-17em} i.e.,\hspace{5em} \quad $-x^2 + 55x - 750 = 0$\\
\vspace{0.4em}
\hspace*{-17em} i.e.,\hspace{6em} \quad $x^2 - 55x + 750 = 0$\\
\vspace{0.4em}
\hspace*{-2em}Therefore, the number of toys produced that day satisfies the quadratic equation
\[
x^2 - 55x + 750 = 0
\]
\hspace{-6em}
which is the required representation of the problem mathematically.

\vspace{0.4em}
\hspace{-6em}
\textbf{\color{cyanblue}Example 2 :} Check whether the following are quadratic equations:
\begin{tabbing}
\hspace{3cm} \= \hspace{6.5cm} \= \kill
\>(i) $(x - 2)^2 + 1 = 2x - 3$ \> (ii) $x(x + 1) + 8 = (x + 2)(x - 2)$ \\
\\
\>(iii) $x(2x + 3) = x^2 + 1$ \> (iv) $(x + 2)^3 = x^3 - 4$
\end{tabbing}
\vspace{0.5em}
\hspace{-28em}
\textbf{\color{cyanblue}Solution :}
\\
\hspace{-10em}(i) LHS = $(x - 2)^2 + 1 = x^2 - 4x + 4 + 1 = x^2 - 4x + 5$\\
\vspace{0.4em}
\hspace*{-11em} Therefore, $(x - 2)^2 + 1 = 2x - 3$ can be rewritten as\\
\vspace{0.4em}
\hspace*{-3em} $x^2 - 4x + 5 = 2x - 3$\\
\vspace{0.4em}
\hspace{-16em} i.e., \hspace{7em}\quad $x^2 - 6x + 8 = 0$\\
\vspace{0.4em}
\hspace{-18em}It is of the form  ax^2 + bx + c = 0.\\
\vspace{0.4em}
\hspace{-10em}Therefore, the given equation is a quadratic equation.\\
\vspace*{5em}
\centering
Reprint 2025-2026


\newpage
\begin{document}
\vspace*{-1cm}
\noindent
{\color{cyanblue}\makebox[\textwidth]{\text{Quadratic Equations} \hfill \text{41}}\\[-0.6em]
\rule{\textwidth}{3pt}}
\noindent\makebox[\textwidth][l]
{\textbf{(ii)} \quad Since \( x(x + 1) + 8 = x^2 + x + 8 \) and \( (x + 2)(x - 2) = x^2 - 4 \) }\\
\vspace{0.4em}
\begin{adjustwidth}{1cm}{0cm}
Therefore,
\hspace{4em}\quad \( x^2 + x + 8 = x^2 - 4 \) \\
\vspace{0.4em}
i.e., 
\hspace{2em}\quad \( x + 12 = 0 \) \\
\vspace{0.4em}
It is not of the form \( ax^2 + bx + c = 0 \). \\
\vspace{0.4em}
Therefore, the given equation is {not} a quadratic equation.
\end{adjustwidth}
\vspace{0.5em}
\noindent\makebox[\textwidth][l]
{\textbf{(iii)} \quad Here,\hspace{6em} \quad LHS \( = x(2x + 3) = 2x^2 + 3x \)} \\
\vspace{-0.4em}
\begin{adjustwidth}{1cm}{0cm}
So, \hspace{7em}\quad \( x(2x + 3) = x^2 + 1 \) can be rewritten as \\
\vspace{0.4em}
\hspace{9em }\quad \( 2x^2 + 3x = x^2 + 1 \) \\
\vspace{0.4em}
Therefore, we get 
\quad \( x^2 + 3x - 1 = 0 \) \\[0.2cm]
It is of the form \( ax^2 + bx + c = 0 \). \\[0.2cm]
So, the given equation is a quadratic equation.
\end{adjustwidth}
\vspace{0.4em}
\noindent\makebox[\textwidth][l]
{\textbf{(iv)} \quad Here, \hspace{6em}\quad LHS \( = (x + 2)^3 = x^3 + 6x^2 + 12x + 8 \)} \\
\begin{adjustwidth}{1cm}{0cm}
Therefore,
\quad \quad \quad \quad \quad  \( (x + 2)^3 = x^3-4 \) can be rewritten as \\[0.2cm]
\hspace{8em}
\quad \( x^3 + 6x^2 + 12x + 8 = x^3 -4 \) \\[0.2cm]
i.e., 
\quad \( 6x^2 + 10x + 6 = 0 \) or,\quad\(x^2+2x+2=0\) \\[0.2cm]
It is of the form \( ax^2 + bx + c = 0 \). \\[0.2cm]
So, the given equation is a quadratic equation.
\end{adjustwidth}\\
{\textcolor{cyanblue}{\textbf{Remark:}} 
{Be careful!} In (ii) above, the given equation appears to be a quadratic equation, but it is not a quadratic equation.
\vspace{-0.4em}
\begin{adjustwidth}{0.5cm}{0cm}
\hspace{2em}
In (iv) above, the given equation appears to be a cubic equation (an equation of degree 3) and not a quadratic equation. But it turns out to be a quadratic equation. As you can see, often we need to simplify the given equation before deciding whether it is quadratic or not.
\end{adjustwidth}
\vspace{1em}
\textcolor{cyanblue}{\textbf{EXERCISE 4.1}}
\vspace{-1em}
\begin{enumerate}
\item Check whether the following are quadratic equations:
\begin{multicols}{}
\begin{enumerate}[label=(\roman*)]
\begin{tabbing}
\hspace{-1cm} \= \hspace{6.5cm} \= \kill
  \>(i) \( x^2 + 1 = 2x \) \>(ii) \( x^2 - 2x = (-2)(3 - x) \)\\
  \>(iii) \( (x - 2)(x + 1) = x(x + 5) \)\>(iv) \( x^2 + (2x - 1)(2x + 3) = x^2 + 2 \)\\
  \>(v) \( x^2 + 1 = (x + 1)^2 \) \>(vi) \( x^3 - 3x^2 + x + 1 = (x - 2)^2 \)\\
  \>(vii) \(x + 2)^3 = 2x (x^2 – 1\)\>(viii)\( x^3– 4x^2 – x + 1 = (x – 2)^3\)
  \end{tabbing}
\end{enumerate}
\end{multicols}
\item Represent the following situations in the form of quadratic equations:
\begin{enumerate}[label=(\roman*)]
\item The area of a rectangular plot is 528 m\(^2\). The length of the plot (in metres) is one more than twice its breadth. We need to find the length and breadth of the plot.
\end{enumerate}
\vspace*{4em}
\centering
Reprint 2025-2026
\end{document}
